% SOURCE_AUTHORITY: sga5_fr_workpass.tex, lines 2311--2451 (LF-normalized unit).
% SOURCE_SHA256: 791F4EFFC5E02832D5D77ED03518C8156D6F07E4C8238B03545DB93D883FBB28
\subsection*{3.1}
Sejam, para $i=1,2$, $X_i$ um $S$-esquema e $L_i\in\ob\Dctf(X_i)$,
\[
X=X_1\times_S X_2,
\qquad
p_i:X\longrightarrow X_i,
\qquad
q_i:X_i\longrightarrow S
\]
as projeções canônicas:
\[
\begin{tikzcd}[column sep=1.8em,row sep=1.5em]
& X \arrow[dl,"P_1"'] \arrow[dr,"P_2"] & \\
L_1\quad X_1 \arrow[dr,"q_1"'] & & X_2\quad L_2 \arrow[dl,"q_2"] \\
& S &
\end{tikzcd}
\]
Com as notações de 1.3, tem-se um isomorfismo canônico (2.2.4)
\[
DL_1\overset{\mathbf L}{\otimes}_S L_2
\xrightarrow{\sim}
\uRHom\bigl(L_1\overset{\mathbf L}{\otimes}_S A_{X_2},
q_1^!A_S\overset{\mathbf L}{\otimes}_S L_2\bigr).
\]
Ao compô-lo com o isomorfismo de Künneth
\[
q_1^!A_S\overset{\mathbf L}{\otimes}_S L_2
\xrightarrow{\sim}p_2^!L_2 \qquad (1.7.3),
\]
obtém-se um isomorfismo
\begin{equation}
DL_1\overset{\mathbf L}{\otimes}_S L_2
\xrightarrow{\sim}
\uRHom(p_1^*L_1,p_2^!L_2).
\tag{3.1.1}
\end{equation}
Esse desempenhará um papel fundamental em tudo o que se segue. Para $X_i=S$, 3.1.1 reduz-se a
\[
L_1^\vee\overset{\mathbf L}{\otimes}L_2\xrightarrow{\sim}\uRHom(L_1,L_2),
\]
caso particular de um isomorfismo válido para complexos perfeitos sobre um topos anelado qualquer (SGA 6 I 7.7). Poremos
\begin{equation}
\uRHom(p_1^*L_1,p_2^!L_2)=\uRHom_S(L_1,L_2),
\qquad
\Hom(p_1^*L_1,p_2^!L_2)=\Hom_S(L_1,L_2).
\tag{3.1.2}
\end{equation}
Os elementos de $\Hom_S(L_1,L_2)$ serão chamados \emph{correspondências cohomológicas de $L_1$ a $L_2$}. Em vez de ``correspondência cohomológica de $A_{X_1}$ a $A_{X_2}$'', diremos simplesmente \emph{correspondência cohomológica de $X_1$ a $X_2$} (subentende-se: com coeficientes em $A$).

Se $X_1$ for \emph{suave}, de dimensão relativa pura $r_1$, o mesmo ocorrerá com $p_2$, e pelo teorema de dualidade (SGA 4 XVIII 3.2.5), tem-se
\[
p_2^!L_2\xrightarrow{\sim}p_2^*L_2(r_1)[2r_1],
\]
portanto
\begin{equation}
\begin{aligned}
&\uRHom_S(L_1,L_2)=\uRHom(p_1^*L_1,p_2^*L_2)(r_1)[2r_1],\\
&\Hom_S(L_1,L_2)=\H^{2r_1}\bigl(X,\uRHom(p_1^*L_1,p_2^*L_2)(r_1)\bigr).
\end{aligned}
\tag{3.1.3}
\end{equation}
Em particular, o grupo das correspondências cohomológicas de $X_1$ a $X_2$ identifica-se neste caso com $\H^{2r_1}(X,A(r_1))$.

\subsection*{3.2}
As correspondências cohomológicas que aparecem na prática estão associadas a $S$-morfismos $X_2\to X_1$, ou, mais geralmente, a ciclos de $X$, ou a esquemas sobre $X$. Seja
\[
c:C\longrightarrow X
\]
um esquema sobre $X$ (às vezes se diz que $c$ é uma correspondência entre $X_1$ e $X_2$), e denotemos por $c_i=p_ic:C\to X_i$ as projeções. Chamaremos \emph{correspondências cohomológicas de $L_1$ a $L_2$ com suporte em $c$} (ou em $C$, se não houver risco de confusão) aos elementos do grupo
\[
\H^0\bigl(X,c_!\uRHom(c_1^*L_1,c_2^!L_2)\bigr)
\quad(=\Hom(c_1^*L_1,c_2^!L_2)\text{ se }c\text{ é próprio}).
\]
Pela fórmula de indução (SGA 4 XVIII 3.1.12.2), tem-se um isomorfismo canônico
\begin{equation}
c^!\uRHom(p_1^*L_1,p_2^!L_2)
\xrightarrow{\sim}
\uRHom(c_1^*L_1,c_2^!L_2),
\tag{3.2.1}
\end{equation}
de onde, ao aplicar $c_!$ e compor com a flecha de adjunção $c_!c^!\to \Id$, obtêm-se flechas canônicas
\begin{equation}
c_!\uRHom(c_1^*L_1,c_2^!L_2)
\longrightarrow
\uRHom(p_1^*L_1,p_2^!L_2),
\tag{3.2.2}
\end{equation}
\begin{equation}
\H^0\bigl(X,c_!\uRHom(c_1^*L_1,c_2^!L_2)\bigr)
\longrightarrow
\Hom_S(L_1,L_2).
\tag{3.2.3}
\end{equation}
Qualquer correspondência de $L_1$ a $L_2$ com suporte em $C$ define, portanto, canonicamente, por meio de 3.2.3, uma correspondência (sem suporte) de $L_1$ a $L_2$.

\emph{Exemplos.} a) Suponhamos que $c$ é próprio e que $c_2$ tem dimensão Tor finita e dimensão relativa $\leq N$. O morfismo de traço (SGA $4^{1/2}$ Ciclo 2.3.3)
\[
\mathrm{Tr}_{c_2}:c_{2!}A_C\longrightarrow A_{X_2}(-N)[-2N]
\]
define por adjunção uma correspondência cohomológica de $A_{X_1}$ a $A_{X_2}(-N)[2N]$ com suporte em $C$, que chamaremos \emph{classe de $(c,N)$}, e denotaremos
\begin{equation}
\operatorname{cl}(c,N)\in \H^{-2N}\bigl(C,c_2^!A_{X_2}(-N)\bigr).
\tag{3.2.4}
\end{equation}
O caso mais importante é aquele em que $N=0$, isto é, $c_2$ é finito (e de dimensão Tor finita), caso em que $\operatorname{cl}(c)$ é uma correspondência cohomológica de $X_1$ a $X_2$ com suporte em $C$. Se $X_1$ for suave, de dimensão relativa pura $r_1$, tem-se, por 3.1.3 e 3.2.1,
\begin{equation}
\H^{-2N}\bigl(C,c_2^!A_{X_2}(-N)\bigr)
=
\H^{2(r_1-N)}\bigl(C,c^!A_X(r_1-N)\bigr).
\tag{3.2.5}
\end{equation}
Se, além disso, $c$ for uma imersão fechada, o segundo membro se reescreverá
\[
\H^{2(r_1-N)}_C\bigl(X,A_X(r_1-N)\bigr),
\]
e a classe de $(c,N)$, considerada como elemento desse grupo, não é senão a classe de cohomologia associada a $c$ no sentido de (SGA $4^{1/2}$ Ciclo 2.3.1), como se deduz imediatamente das definições.

b) Seja $F:X_2\to X_1$ um $S$-morfismo, e denotemos por $c(F):C(F)\to X$ seu gráfico, isto é, a imagem da seção de $p_2$ definida por $F$ (``o conjunto dos pontos $(Fx,x)$''): $c_2$ é um isomorfismo, e $c_1c_2^{-1}=F$. Têm-se, portanto, isomorfismos canônicos
\begin{equation}
\uRHom\bigl(c_1(F)^*L_1,c_2(F)^!L_2\bigr)=\uRHom(F^*L_1,L_2),
\qquad
\Hom\bigl(c_1(F)^*L_1,c_2(F)^!L_2\bigr)=\Hom(F^*L_1,L_2).
\tag{3.2.6}
\end{equation}
Em outros termos, as correspondências cohomológicas de $L_1$ a $L_2$ com suporte no gráfico de $F$ são os ``levantamentos'' de $F$ a um homomorfismo $F^*L_1\to L_2$. Eis dois exemplos importantes de tais correspondências:

(i) $X_1=X_2$, $L_1=L_2$, $F=1$, a identidade de $X_1$, cujo gráfico $\Delta$ é a diagonal. A correspondência com suporte em $\Delta$ definida pela identidade de $L_1$ chama-se \emph{correspondência diagonal}. É o elemento
\[
1\in \Hom(L_1,L_1)=\H^0\bigl(X,\Delta_*\Delta^!\uRHom_S(L_1,L_1)\bigr)
\]
(grupo que se identifica com $\H^{2r}_\Delta(X,\uRHom_S(L_1,L_1)(r))$ quando $X_1$ é suave, de dimensão relativa pura $r$).

(ii) Suponhamos que $S$ seja o espectro de um corpo finito $\mathbb F_q$, que $X_1=X_2$, $L_1=L_2$, e tomemos como $F:X_1\to X_1$ o endomorfismo de Frobenius relativo a $\mathbb F_q$, identidade sobre o espaço subjacente e elevação à potência $q$-ésima sobre $\mathcal O_{X_1}$ (cf. XIV), e como levantamento de $F$ a $L_1$ o isomorfismo canônico
\[
F^*L_1\xrightarrow{\sim}L_1
\]
(loc. cit. §2 prop. 1) (mais precisamente, o iterado $a$-ésimo do isomorfismo de (loc. cit.), se $q=p^a$). A correspondência cohomológica de $L_1$ a $L_1$ com suporte no gráfico de $F$ assim definida chama-se \emph{correspondência de Frobenius} (relativa a $\mathbb F_q$).

c) Algumas correspondências cohomológicas entre feixes sobre curvas, relacionadas aos operadores de Hecke, foram estudadas por Langlands em ([13] §7).

d) As correspondências cohomológicas entre feixes constantes (as mais frequentes na natureza) fornecem naturalmente, por ``integração'', correspondências cohomológicas entre complexos, como veremos agora.
